\documentclass[exercice,a4]{thedocument}%
\usepackage{texmaths-tikz}%
\usepackage{pstricks}%
\usepackage{pst-bezier}%
\usepackage{pstricks-add}%
\usepackage{pst-eucl}%
\usepackage{pst-func}%
\usepackage{pst-tree}%
\usepackage{pst-circ}%
\usepackage{tikz}%
\usepackage{pgfplots}%
\usepackage{tkz-tab}%
\usepackage{tikzpagenodes}%
\usepackage{tabularray}%
\usepackage{tabularx}%
\usepackage{multirow}%
\usepackage{multicol}%
\usepackage{scratch3}%
\usepackage{enumitem}%
\usepackage[np]{numprint}%
\usepackage{pas-tableur}%
\usepackage{listings}%
\usepackage{graphicx}%
\usepackage{adjustbox}%
\usepackage{needspace}%
\usepackage{xcolor}%
\usepackage{eurosym}%
\usepackage{diagbox}%
\usepackage{colortbl}%
\usepackage{multido}%
\usepackage{icomma}%
\usepackage[french]{babel}%
\startdatakeys
    \source{26GENMATAN1}
    \BO{2026}
    \cycle{4}
    \competencesDuSocle{CA,RA,CO}
    \objectifsApprentissage{C4A1aa,C4A2ad,C4B2cb,C4D2aa,C4D3ac}
    \motsCles{Scratch,programme,carré,triangle équilatéral,angle}
\enddatakeys
\definecolor{gris}{gray}{0.5}%
\providecommand{\wideparen}[1]{\overset{\frown}{#1}}%
\begin{document}%
\providecommand{\degres}{^\circ}%
Dans cet exercice, aucune justification n'est demandée.

Un élève souhaite réaliser une figure constituée de carrés et de triangles équilatéraux, à l'aide d'un logiciel de programmation. Pour cela, il crée les trois blocs ci-dessous :

\vspace{10pt}

\begin{center}
\begin{tabularx}{\linewidth}{|*{3}{X|}}\hline
\multicolumn{1}{|c|}{\textbf{Bloc 1}} &
\multicolumn{1}{|c|}{\textbf{Bloc 2}} &\multicolumn{1}{|c|}{\textbf{Bloc 3}}\\ \hline
\smallskip
\begin{scratch}[scale=0.86]
\initmoreblocks{définir \namemoreblocks{Initialisation}}
\blockmove{aller à x : \ovalnum{0} y : \ovalnum{0}}
\blockmove{s'orienter à ~\ovalnum{90}}
\blockpen{effacer tout}
\blockpen{stylo en position d'écriture}
\end{scratch}
&
\smallskip
\begin{scratch}[scale=0.86]
\initmoreblocks{définir \namemoreblocks{carré}}
\blockrepeat{répéter \ovalnum{\textbf{A}} fois}
{
\blockmove{avancer de \ovalnum{50} pas}
\blockmove{tourner \turnleft{} de \ovalnum{\textbf{B}} degrés}
}
\end{scratch}
&
\smallskip
\begin{scratch}[scale=0.86]
\initmoreblocks{définir \namemoreblocks{triangle équilatéral}}
\blockrepeat{répéter \ovalnum{\textbf{C}} fois}
{
\blockmove{avancer de \ovalnum{50} pas}
\blockmove{tourner \turnright{} de \ovalnum{\textbf{D}} degrés}
}
\end{scratch}
\\ \hline
\end{tabularx}
\end{center}

\medskip


L'instruction \og s'orienter à 90 \fg{} signifie que le lutin se dirige vers la droite.

\begin{enumerate}
\item Quelles sont les coordonnées du lutin après l'exécution du Bloc 1\,?
\item Dans les blocs 2 et 3, on a remplacé certaines valeurs par les lettres \textbf{A}, \textbf{B}, \textbf{C} et \textbf{D}.

Sur la copie, indiquer la lettre et sa valeur correspondante.
\item L'élève a construit trois figures avec les trois programmes ci-dessous.

Associer chaque figure au programme correspondant.
\end{enumerate}

\medskip

\begin{center}
\begin{tabularx}{\linewidth}{|*{3}{X|}}\hline
\multicolumn{1}{|c|}{\textbf{Programme 1}} &
\multicolumn{1}{|c|}{\textbf{Programme 2}} &\multicolumn{1}{|c|}{\textbf{Programme 3}}\\ \hline
\begin{scratch}[scale=0.9]
\blockinit{quand \greenflag est cliqué}
\blockmoreblocks{Initialisation}
\blockmoreblocks{triangle équilatéral}
\blockrepeat{répéter \ovalnum{\textbf{3}} fois}
{
\blockmoreblocks{carré}
\blockmove{avancer de \ovalnum{50} pas}
\blockmove{tourner \turnright{} de \ovalnum{\textbf{120}} degrés}
}
\end{scratch}
&
\begin{scratch}[scale=0.9]
\blockinit{quand \greenflag est cliqué}
\blockmoreblocks{Initialisation}
\blockmoreblocks{carré}
\blockrepeat{répéter \ovalnum{\textbf{4}} fois}
{
\blockmoreblocks{triangle équilatéral}
\blockmove{avancer de \ovalnum{50} pas}
\blockmove{\turnleft{} de \ovalnum{\textbf{90}} degrés}
}
\end{scratch}
&
\begin{scratch}[scale=0.9]
\blockinit{quand \greenflag est cliqué}
\blockmoreblocks{Initialisation}
\blockmoreblocks{triangle équilatéral}
\blockrepeat{répéter \ovalnum{\textbf{3}} fois}
{
\blockmove{avancer de \ovalnum{50} pas}
\blockmove{tourner \turnright{} \ovalnum{\textbf{60}} degrés}
\blockmoreblocks{triangle équilatéral}
\blockmove{tourner \turnright{} \ovalnum{\textbf{60}} degrés}
}
\end{scratch}\\ \hline
\end{tabularx}

\vspace{14pt}

\begin{tabularx}{\linewidth}{|*{3}{>{\centering \arraybackslash}X|}}\hline
\textbf{Figure A} & \textbf{Figure B} &\textbf{Figure C}\\ \hline
\begin{center}
%%début Tikz
%\begin{tikzpicture}[scale=0.4]
%% Triangle du bas
%\draw[thick] (0,0) -- (2,0) -- (1,1.732) -- cycle;
%% Triangle gauche
%\draw[thick] (0,0) -- (-2,0) -- (-1,1.732) -- cycle;
%% Triangle droit
%\draw[thick] (2,0) -- (4,0) -- (3,1.732) -- cycle;
%% Triangle du bas central inversé (ligne commune)
%% Petit triangle central pointe en bas
%\draw[thick] (0,1.732) -- (2,1.732) -- (1,0) -- cycle;
%\end{tikzpicture}
%Fin Tikz
%PSTricks
\psset{unit=1cm}
\begin{pspicture}(-1.5,-1.5)(1.5,1.5)
\pspolygon(1;-30)(1;90)(1;210)
\pspolygon(0.5;30)(0.5;150)(0.5;270)
\end{pspicture}
%fin PSTricks
\end{center}
\vspace{4pt}
&
\vspace{6pt}
\begin{center}
%%début Tikz
%\begin{tikzpicture}[scale=0.4]
%% Figure B : carré + triangles sur chaque côté, répété 4 fois (style losange)
%\draw[thick] (0,0) -- (2,0) -- (2,2) -- (0,2) -- cycle;
%\draw[thick] (0,0) -- (-1,-1.732) -- (1,-1.732) -- (2,0);
%\draw[thick] (2,0) -- (3.732,1) -- (2.732,2.732) -- (2,2);
%\end{tikzpicture}
%Fin Tikz
%PSTricks
\psset{unit=1cm}
\begin{pspicture}(-1.5,-1.5)(1.5,1.5)
%\pspolygon(1;-30)(1;90)(1;210)
\pspolygon(0.45;30)(0.45;150)(0.45;270)
\def\car{\psframe(-0.42,-0.42)(0.42,0.4)}
\multido{\n=90+120}{3}{\rput{\n}(0.66;\n){\car}}
\end{pspicture}
%fin PSTricks
\end{center}
\vspace{4pt}
&
\vspace{6pt}
\begin{center}
%%début Tikz
%\begin{tikzpicture}[scale=0.4]
%% Figure C : deux triangles se croisant (étoile)
%\draw[thick] (-1.5,0) -- (1.5,0) -- (0,2.6) -- cycle;
%\draw[thick] (-1.5,1.3) -- (1.5,1.3) -- (0,-1.3) -- cycle;
%\end{tikzpicture}
%Fin Tikz
%PSTricks
\psset{unit=1cm}
\begin{pspicture}(-1.5,-1.5)(1.5,1.5)
%\pspolygon(1;-30)(1;90)(1;210)
\def\tri{\pspolygon(0.57;-30)(0.57;90)(0.57;210)}
%\def\car{\psframe(-0.42,-0.42)(0.42,0.4)}
\psframe(-0.49,-0.49)(0.49,0.49)
%\multido{\n=0+90}{4}{\rput{\n}(0.66;\n){\tri}}
\rput(0,0.768){\tri}\rput{90}(-0.768,0){\tri}
\rput{180}(0,-0.768){\tri}\rput{-90}(0.768,0){\tri}
\end{pspicture}
%fin PSTricks
\end{center}
\vspace{6pt}\\ \hline
\end{tabularx}
\end{center}
\medskip
\showthesource
\end{document}
